% Part: incompleteness
% Chapter: arithmetization-syntax
% Section: proofs-in-nd

\documentclass[../../../include/open-logic-section]{subfiles}

\begin{document}

\olfileid[tr]{inc}{art}{pnd}
\olsection{Doğal Tümdengelimde \usetoken{P}{derivation}}

\begin{explain}
!!{derivation}s{}i aritmetikleştirmek için !!{derivation}s{}i sayılarla temsil
etmemiz gerekir. !!{derivation}s, her çıkarımın bir ya da iki etiket taşıdığı
!!{formula}s ağaçları olduğundan, en açık yaklaşım özyinelemeli bir temsildir:
Bir !!{derivation}{}i, bileşenleri son çıkarımın öncüllerine götüren dolaysız
alt-!!{derivation}s{}in sayısı, bu alt-!!{derivation}s{}in temsilleri, son
!!{formula}, son çıkarımın kapatma etiketi ve son çıkarımın türünü belirten
bir sayı olan bir demetle temsil ederiz.
\end{explain}

\begin{defn}
  $\delta$, doğal tümdengelimde !!a{derivation} ise $\Gn{\delta}$ tümevarımlı
  olarak aşağıdaki gibi tanımlanır:
  \begin{enumerate}
  \item
    $\delta$ yalnızca $!A$ varsayımından oluşuyorsa $\Gn{\delta}$
    $\tuple{0, \Gn{!A}, n}$ olur. $n$ sayısı, varsayım !!{undischarged} ise
    $0$, değilse varsayımın sayısal etiketidir.
  \item
    $\delta$ sıfır, bir, iki ya da üç öncüllü bir çıkarımla bitiyorsa
    $\Gn{\delta}$ sırasıyla
    \begin{align*}
      & \tuple{0, \Gn{!A}, n, k}, \\
      & \tuple{1, \Gn{\delta_1}, \Gn{!A}, n, k}, \\
      & \tuple{2, \Gn{\delta_1}, \Gn{\delta_2}, \Gn{!A}, n, k}, \text{ ya da}\\
      & \tuple{3, \Gn{\delta_1}, \Gn{\delta_2}, \Gn{\delta_3}, \Gn{!A}, n,
        k}
    \end{align*}
    olur. Burada $\delta_1$, $\delta_2$, $\delta_3$, $\delta$'daki son
    çıkarımın öncülüyle ya da öncülleriyle biten alt-!!{derivation}s;
    $!A$, $\delta$'daki son çıkarımın sonucu; $n$, son çıkarımın kapatma
    etiketi (çıkarım hiçbir varsayımı kapatmıyorsa $0$); $k$ ise son çıkarımda
    kullanılan kurala göre aşağıdaki tabloyla verilen sayıdır.

    \begin{tabular}{lccccccc}
      \text{Kural:} & \Intro{\land} & \Elim{\land} & \Intro{\lor} & \Elim{\lor} \\
      $k$: & 1 & 2 & 3 & 4  \\[1.5ex]
      \text{Kural:} &  \Intro{\lif} & \Elim{\lif} & \Intro{\lnot} & \Elim{\lnot} \\
      $k$: & 5 & 6 & 7 & 8 \\[1.5ex]
      \text{Kural:}  &  \FalseInt & \FalseCl & \Intro{\lforall} & \Elim{\lforall} \\
      $k$: & 9 & 10 & 11 & 12 \\[1.5ex]
      \text{Kural:} & \Intro{\lexists} & \Elim{\lexists} & \Intro{\eq} & \Elim{\eq} \\
      $k$: & 13 & 14 & 15 & 16
    \end{tabular}
  \end{enumerate}
\end{defn}

\begin{ex}
  Şu çok basit !!{derivation}{}i ele alın:
  \begin{prooftree}
    \AxiomC{$\Discharge{!A \land !B}{1}$}
    \RightLabel{\Elim{\land}}
    \UnaryInfC{$!A$}
    \DischargeRule{\Intro{\lif}}{1}
    \UnaryInfC{$(!A \land !B) \lif !A$}
  \end{prooftree}
  Varsayımın Gödel numarası $d_0 = \tuple{0,
    \Gn{!A \land !B}, 1}$ olurdu. $\Elim{\land}$ sonucuyla biten
  !!{derivation}{}in Gödel numarası $d_1 = \tuple{1,
     d_0, \Gn{!A}, 0, 2}$ olurdu ($\Elim{\land}$ bir öncüllü olduğu için
  $1$, sonuç $!A$'nın Gödel numarası, hiçbir varsayım kapatılmadığı için
  $0$ ve $\Elim{\land}$'ı kodlayan sayı olduğu için $2$). Böylece bütün
  !!{derivation}{}in Gödel numarası
  $\tuple{1, d_1, \Gn{((!A \land !B) \lif
      !A)}, 1, 5}$, yani
  \[
  \tuple{1, \tuple{1, \tuple{0, \Gn{(!A \land !B)}, 1}, \Gn{!A}, 0, 2},
    \Gn{((!A \land !B) \lif !A)}, 1, 5}
  \]
  olur.
\end{ex}

\begin{explain}
!!{derivation}s için bir temsil belirlediğimize göre bu !!{derivation}s{}in
Gödel numaralarını ilkel özyinelemeli biçimde işleyebildiğimizi ve temel
özellikleriyle bağıntılarını ifade edebildiğimizi de göstermeliyiz. Bazı
işlemler basittir: Örneğin !!a{derivation}{}in $d$ Gödel numarası verildiğinde
$\fn{EndFmla}(d) = (d)_{(d)_0+1}$ bize son !!{formula} için Gödel numarasını,
$\fn{DischargeLabel}(d) = (d)_{(d)_0 + 2}$ kapatma etiketini ve
$\fn{LastRule}(d) = (d)_{(d)_0 + 3}$ son çıkarımın türünü belirten sayıyı
verir. Bazı işlemler ise çok daha zordur. En azından bunların nasıl
yapılacağını ana hatlarıyla göstereceğiz. Amaç, ``$\delta$, $\Gamma$'dan
$!A$'nın !!a{derivation}{}idir'' bağıntısının $\delta$ ile $!A$'nın Gödel
numaraları üzerinde ilkel özyinelemeli bir bağıntı olduğunu göstermektir.
\end{explain}

\begin{prop}
  Aşağıdaki bağıntılar ilkel özyinelemelidir:
  \begin{enumerate}
  \item $!A$, $\delta$ içinde $n$ etiketiyle bir varsayım olarak geçer.
  \item $\delta$ içindeki $n$ etiketli bütün varsayımlar $!A$ biçimindedir
    (yani $\delta$ içinde $n$ etiketini kullanarak $!A$ varsayımını
    !!{discharge} yapabiliriz).
  \end{enumerate}
\end{prop}

\begin{proof}
  !!{formula}s{}in Gödel numaraları ile !!{derivation}s{}in Gödel numaraları
  arasındaki karşılık gelen bağıntıların ilkel özyinelemeli olduğunu
  göstermeliyiz.
  \begin{enumerate}
  \item Gödel numarası $x$ olan bir formülün, Gödel numarası $d$ olan
    !!{derivation}{}in $n$ etiketli bir varsayımı olması durumunda geçerli olan
    $\fn{Assum}(x, d, n)$ bağıntısının ilkel özyinelemeli olduğunu göstermek
    istiyoruz. Bu, Gödel numarası $\tuple{0, x, n}$ olan !!{derivation},
    $d$'nin bir alt-!!{derivation}{}i olduğunda geçerlidir. !!{derivation}s{}i
    kodlama biçimimizin \olref[cmp][rec][tre]{sec}'te tanıtılan ağaç
    kodlamasının özel bir durumu olduğuna dikkat edin. Dolayısıyla ilkel
    özyinelemeli $\fn{SubtreeSeq}(d)$ fonksiyonu, $d$'nin bütün
    alt-!!{derivation}s{}inin (uzunluğu en çok $d$ olan) Gödel numaraları
    dizisini verir. Böylece
    \[
    \fn{Assum}(x, d, n) \defiff \bexists{i<d}{(\fn{SubtreeSeq}(d))_i =
      \tuple{0, x, n}}
    \]
    olarak tanımlayabiliriz.
  \item Gödel numarası $d$ olan !!{derivation} içindeki $n$ etiketli bütün
    varsayımların Gödel numarası $x$ olan !!{formula} olması durumunda geçerli
    olan $\fn{Discharge}(x, d, n)$ bağıntısının ilkel özyinelemeli olduğunu
    göstermek istiyoruz. Bu bağıntı, ancak ve ancak
    $\bforall{y<d}{(\fn{Assum}(y, d, n) \lif y = x)}$ geçerliyse geçerlidir.
  \end{enumerate}
\end{proof}

\begin{prop}
  \ollabel{prop:followsby}
  Gödel numarası $d$ olan $\delta$ !!{derivation}{}indeki son çıkarımın doğru
  olması durumunda geçerli olan $\fn{Correct}(d)$ özelliği ilkel
  özyinelemelidir.
\end{prop}

\begin{proof}
  Burada her $R$ çıkarım kuralı için $\fn{FollowsBy}_R(d)$ bağıntısının
  ilkel özyinelemeli olduğunu göstermeliyiz. $\fn{FollowsBy}_R(d)$, ancak ve
  ancak $d$, bir $\delta$ !!{derivation}{}inin Gödel numarasıysa ve
  $\delta$'nın sonundaki !!{formula}, $\delta$'nın dolaysız
  alt-!!{derivation}s{}inden $R$'nin doğru uygulanmasıyla çıkıyorsa geçerlidir.

  Basit bir durum $\Intro{\land}$ kuralıdır. $\delta$ doğru bir
  $\Intro{\land}$ çıkarımıyla bitiyorsa şu biçimdedir:
  \begin{prooftree}
    \AxiomC{}
    \RightLabel{$\delta_1$}
    \DeduceC{$!A$}

    \AxiomC{}
    \RightLabel{$\delta_2$}
    \DeduceC{$!B$}

    \RightLabel{\Intro\land}
    \BinaryInfC{$!A \land !B$}
  \end{prooftree}
  Bu durumda $\delta$'nın Gödel numarası $d = \tuple{2, d_1, d_2,
    \Gn{(!A \land !B)}, 0, k}$ olur; burada
  $\fn{EndFmla}(d_1) = \Gn{!A}$, $\fn{EndFmla}(d_2) = \Gn{!B}$,
  $n=0$ ve $k=1$'dir. Dolayısıyla
  $\fn{FollowsBy}_{\Intro\land}(d)$ bağıntısını
  \begin{multline*}
    (d)_0 = 2 \land \fn{DischargeLabel}(d) = 0 \land \fn{LastRule}(d) = 1 \land {}\\
    \fn{EndFmla}(d) = {}\\
    \Gn{(} \concat \fn{EndFmla}((d)_1) \concat \Gn{\land}
    \concat \fn{EndFmla}((d)_2) \concat \Gn{)}
  \end{multline*}
  olarak tanımlayabiliriz.

  Bir başka basit örnek $\Intro\eq$ kuralıdır. Bunun öncülü yoktur; dolayısıyla
  varsayımlarda olduğu gibi $(d)_0 = 0$ olur. Ayrıca kapatma etiketi de
  yoktur, yani $n=0$'dır. Ancak $!A$, kapalı bir $t$ terimi için
  $\eq[t][t]$ biçiminde olmalıdır. İlkel özyinelemeli bir tanım şöyledir:
  \begin{multline*}
    (d)_0 = 0 \land \fn{DischargeLabel}(d) = 0 \land {}\\
    \bexists{t<d}{(\fn{ClTerm}(t) \land
      \fn{EndFmla}(d) = {}\\
      \Gn{{\eq}(} \concat t \concat \Gn{,} \concat t \concat \Gn{)})}
  \end{multline*}

  Daha karmaşık bir örnek için $\fn{FollowsBy}_{\Intro{\lif}}(d)$, ancak ve
  ancak $\delta$'nın sonundaki !!{formula} $(!A \lif !B)$ biçimindeyse,
  $\delta_1$'in sonundaki !!{formula} $!B$ ise ve $\delta$'da $n$ ile etiketli
  her varsayım $!A$ biçimindeyse geçerlidir. Bunu ilkel özyinelemeli biçimde
  şöyle ifade edebiliriz:
  \begin{multline*}
    (d)_0 = 1 \land {}\\
    \bexists{a<d}{(\fn{Discharge}(a, (d)_1, \fn{DischargeLabel}(d)) \land {}}\\
      \fn{EndFmla}(d) = (\Gn{(} \concat a \concat \Gn{\lif}
      \concat \fn{EndFmla}((d)_1) \concat \Gn{)}))
  \end{multline*}
  (Burada $a$'yı $!A$'nın Gödel numarası olarak düşünün.)

  Bir başka örnek olarak \Intro{\lexists} kuralını ele alın. Burada
  $\delta$'daki son çıkarım, ancak ve ancak $\Subst{!A}{t}{x}$,
  $\delta_1$ !!{derivation}{}inin sonunda bulunan !!{formula} ve
  $\lexists[x][!A]$ son çıkarımın sonucu olacak biçimde !!a{formula}~$!A$,
  kapalı bir $t$ terimi ve !!a{variable}~$x$ varsa doğrudur. Dolayısıyla
  $\fn{FollowsBy}_{\Intro{\lexists}}(d)$ ancak ve ancak
  \begin{multline*}
    (d)_0 = 1 \land \fn{DischargeLabel}(d) = 0 \land {} \\
    \bexists{a < d}{\bexists{x<d}{\bexists{t<d}{
          (\fn{ClTerm}(t) \land \fn{Var}(x)  \land {}}}}\\
    \fn{Subst}(a,t,x) = \fn{EndFmla}((d)_1) \land
    \fn{EndFmla}(d) = (\Gn{\lexists} \concat x \concat a))
  \end{multline*}
  geçerliyse geçerlidir.

  Ardından $\fn{Correct}(d)$'yi şöyle tanımlarız:
  \begin{multline*}
    \fn{Sent}(\fn{EndFmla}(d)) \land {}\\
    [(\fn{LastRule}(d) = 1 \land
    \fn{FollowsBy}_{\Intro\land}(d)) \lor \dots \lor {}\\
    (\fn{LastRule}(d) = 16 \land \fn{FollowsBy}_{\Elim\eq}(d)) \lor {}\\
    \bexists{n<d}{\bexists{x<d}{(d = \tuple{0, x, n})}}]
  \end{multline*}
  İlk satır $d$'nin sonundaki !!{formula} ifadesinin bir tümce olmasını sağlar. Son satır
  ise $d$'nin yalnızca bir varsayım olduğu durumu kapsar.
\end{proof}

\begin{prob}
  Aşağıdaki özellikleri \olref[inc][art][pnd]{prop:followsby}'deki gibi
  tanımlayın:
  \begin{enumerate}
  \item $\fn{FollowsBy}_{\Elim{\lif}}(d)$,
  \item $\fn{FollowsBy}_{\Elim{\eq}}(d)$,
  \item $\fn{FollowsBy}_{\Elim{\lor}}(d)$,
  \item $\fn{FollowsBy}_{\Intro{\lforall}}(d)$.
  \end{enumerate}
  Sonuncusu için Gödel numarası $d$ olan !!{derivation}{}in son çıkarımının
  özdeğişken koşulunu sağlayıp sağlamadığını, yani
  $\Intro{\lforall}$ çıkarımının $a$ özdeğişkeninin ne $d$'nin sonundaki
  !!{formula} içinde ne de $d$'nin açık bir varsayımında geçtiğini ilkel
  özyinelemeli biçimde sınayabildiğinizi de göstermelisiniz. Bunun için
  \olref[inc][art][pnd]{prop:openassum}'daki ilkel özyinelemeli
  $\fn{OpenAssum}$ yüklemini kullanabilirsiniz.
\end{prob}

\begin{prop}
  \ollabel{prop:deriv}
  $d$'nin doğru bir $\delta$ !!{derivation}{}inin Gödel numarası olması
  durumunda geçerli olan $\fn{Deriv}(d)$ bağıntısı ilkel özyinelemelidir.
\end{prop}

\begin{proof}
  !!^a{derivation}~$\delta$, ancak ve ancak çıkarımlarının her biri bir
  kuralın doğru uygulamasıysa, başka bir deyişle alt-!!{derivation}s{}inin her
  biri doğru bir çıkarımla bitiyorsa doğrudur. Dolayısıyla $\fn{Deriv}(d)$
  ancak ve ancak
  \[
  \bforall{i<\len{\fn{SubtreeSeq}(d)}}{\fn{Correct}((\fn{SubtreeSeq}(d))_i)}
  \]
  geçerliyse geçerlidir.
\end{proof}

\begin{prop}
  \ollabel{prop:openassum} Gödel numarası $z$ olan bir $!A$ varsayımının,
  Gödel numarası $d$ olan $\delta$ !!{derivation}{}inin !!a{undischarged}
  varsayımı olması durumunda geçerli olan $\fn{OpenAssum}(z, d)$ bağıntısı
  ilkel özyinelemelidir.
\end{prop}

\begin{proof}
  Bir varsayım geçişi, $\delta$'nın kapatma etiketi $n$ olan bir kuralla
  biten alt-!!{derivation}{}inde $n$ etiketiyle geçiyorsa !!{discharged}{}tır.
  Dolayısıyla $!A$, ancak ve ancak geçişlerinden en az biri $\delta$ içinde
  !!{discharged} değilse $\delta$'nın !!a{undischarged} varsayımıdır. Dikkatli
  olmalıyız: $\delta$, $!A$'nın hem !!{discharged} hem de !!{undischarged}
  geçişlerini içerebilir.

  $\delta_0 = \delta$ ve $\delta_k$, bir $n$ için
  $\Discharge{!A}{n}$ varsayımı olacak ve $\delta_{i+1}$,
  $\delta_i$'nin dolaysız bir alt-!!{derivation}{}i olacak biçimde bir
  $\delta_0$, \dots, $\delta_k$ dizisini ele alın. Hiçbir $\delta_i$'nin
  kapatma etiketi $n$ olan bir çıkarımla bitmediği böyle bir dizi varsa
  $!A$, $\delta$'nın !!a{undischarged} varsayımıdır.

  İlkel özyinelemeli $\fn{SubtreeSeq}(d)$ fonksiyonu bize $\delta$'nın bütün
  alt-!!{derivation}s{}inin Gödel numaraları dizisini verir. $\delta$'nın
  alt-!!{derivation}s{}inin Gödel numaralarından oluşan her dizi bunun bir alt
  dizisidir. Alt dizi olma, ilkel özyinelemeli bir bağıntıdır:
  $\fn{Subseq}(s, s')$ ancak ve ancak
  $\bforall{i<\len{s}}{\lexists[j<\len{s'}][(s)_i = (s')_j]}$ geçerliyse
  geçerlidir. Dolaysız alt-!!{derivation} olma da öyledir:
  $\fn{Subderiv}(d, d')$ ancak ve ancak
  $\bexists{j<(d')_0}{d = (d')_{j+1}}$ geçerliyse geçerlidir. Dolayısıyla
  $\fn{OpenAssum}(z, d)$'yi şöyle tanımlayabiliriz:
  \begin{multline*}
    \bexists{s<\fn{SubtreeSeq}(d)}{(\fn{Subseq}(s, \fn{SubtreeSeq}(d))
      \land (s)_0 = d \land {}} \\
    \bexists{n<d}{((s)_{\len{s} \tsub 1} = \tuple{0, z, n} \land {}}\\
      \bforall{i<(\len{s} \tsub 1)}{(\fn{Subderiv}((s)_{i+1}, (s)_i) \land {}}\\
      \fn{DischargeLabel}((s)_i) \neq n)))
  \end{multline*}
\end{proof}

\begin{prop}
  \ollabel{prop:prf-prim-rec}
  $\Gamma$'nın ilkel özyinelemeli bir !!{sentence}s kümesi olduğunu varsayın.
  O zaman ``$x$, $!A$'nın !!a{derivation}{}i olan ve !!{undischarged}
  varsayımları $\Gamma$ içinde bulunan bir $\delta$'nın kodudur ve $y$,
  $!A$'nın Gödel
  numarasıdır'' ifadesini veren $\Prf[\Gamma](x, y)$ bağıntısı ilkel
  özyinelemelidir.
\end{prop}

\begin{proof}
  ``$y \in \Gamma$'' ifadesinin ilkel özyinelemeli $R_\Gamma(y)$ yüklemiyle
  verildiğini varsayın. Ancak ve ancak $y$, bir $!A$ tümcesinin Gödel
  numarasıysa ve $x$, son !!{formula} olarak $!A$'yı taşıyan, bütün !!{undischarged}
  varsayımları da $\Gamma$'da bulunan bir doğal tümdengelim !!{derivation}{}inin
  koduysa geçerli olan $\Prf[\Gamma](x, y)$ bağıntısının ilkel
  özyinelemeli olduğunu göstermeliyiz.

  \olref{prop:deriv}'e göre $x$'in doğal tümdengelimde doğru bir $\delta$
  !!{derivation}{}inin Gödel numarası olması durumunda geçerli olan
  $\fn{Deriv}(x)$ özelliği ilkel özyinelemelidir. Dolayısıyla
  $\Prf[\Gamma](x, y)$'yi şöyle tanımlayabiliriz:
  \begin{align*}
    \Prf[\Gamma](x, y) \defiff {}
    & \fn{Deriv}(x) \land \fn{EndFmla}(x) = y \land {} \\
    & \bforall{z < x}{(\fn{OpenAssum}(z, x) \lif R_\Gamma(z))}
  \end{align*}
\end{proof}

\end{document}
